%This study investigated whether mobile network cells can be pre-classified into recreational and urban groups based on traffic behaviour alone, under the assumption that recreational cells are seasonal-dominant and urban cells are non-seasonal. It also examined whether this pre-classification improves forecasting performance. The results show that meaningful behavioural differences exist and that cells can be reliably separated using traffic data alone. However, the forecasting results indicate that the benefit of segmentation is limited compared to the impact of feature engineering and model selection. \\

The t-test approach (Figure \ref{fig:t_test_all_cells}) successfully separated 88 out of 91 cells into seasonal and non-seasonal groups using differences between holiday and non-holiday traffic. The resulting labels were evaluated against geographic categories used as a proxy reference. This suggests that holiday sensitivity captures a meaningful behavioural characteristic of network traffic. At the same time, several exceptions highlight the limitations of relying on geographic labels for validation. Cells such as R3 and U3 exhibited traffic patterns that differed from their expected geographic category, indicating that location alone may not fully describe how a cell is used. In this sense, the behavioural classification may capture operational characteristics that are not visible through geographic categorisation. \\

The forecasting results show that feature-based statistical and machine learning models outperform naive baseline approaches. This demonstrates that the dataset contains strong and predictable temporal structure that can be captured through engineered features. \\

A key finding is the relatively small performance difference between OLS and XGBoost. While XGBoost achieves the lowest overall RMSE and outperforms OLS on the non-seasonal cells (1.56 vs 1.60), the improvements remains modest. On seasonal cells, the paired t-test indicates that there is no statistical difference between OLS and XGBoost with a p-value of 0.42. This suggests that most of the predictive performance can be captured by a simple linear regression with strong feature engineering. Lagged observations and rolling statistics appear to contain most of the relevant information, allowing a linear model to achieve performance comparable to more complex approaches. \\

The consistently weaker performance of SARIMAX further supports this interpretation. While SARIMAX incorporates external regressors, its core structure still relies on autoregressive and moving-average components to model temporal dependence. In contrast, OLS and XGBoost directly leverage a set of engineered features, including lagged observations, rolling statistics and calender effects. The results suggest that these engineered variables capture much of the relevant temporal structure more effectively than the internal ARIMA dynamics.These features appear to provide a more direct representation of temporal dependence than relying on autoregressive structure within the model itself. \\

The impact of dataset segmentation is limited. Training separate models for seasonal and non-seasonal cells in the statistical models only results in marginal improvements compared to a global model, and only for the seasonal subset, which represents the smaller minority group. When a subgroup is already well represented in the global training data, as is the case for non-seasonal cells at 69\% of all samples, group-specific training provides little to no additional benefit. This suggests that pre-classification is most valuable when a meaningful class imbalance exists, which does not appear to be the case in this dataset. \\

Overall, the results show that a simple pipeline combining statistical classification and feature-driven forecasting models, whether linear or tree-based, is sufficient to achieve strong performance in this setting when comparing to baseline models. Feature engineering plays a more important role than model complexity or segmentation of the data, although model choice still produces small, statistically significant differences depending on the seasonality of the subgroup. \\




